\documentclass{article}
\usepackage{iclr2026_conference,times}
\input{math_commands.tex}
\usepackage{hyperref}
\usepackage{url}
\usepackage{graphicx}
\usepackage{booktabs}
\usepackage{amsmath,amssymb}
\usepackage{xcolor}

\title{SCDE: Statistics-Conditioned Decoding for\\ Strict Zero-Shot Cross-Subject EEG Emotion Recognition}

\author{Anonymous authors\\ Paper under double-blind review}

\newcommand{\method}{SCDE}
\newcommand{\scn}{SCN}

\begin{document}

\maketitle

\begin{abstract}
Cross-subject EEG emotion recognition is dominated by an uncomfortable
practical gap: methods that report high accuracies routinely rely on target
subjects' unlabeled data for calibration, while the strict zero-shot setting
--- where the test subject contributes \emph{nothing}, not even unlabeled
signal statistics --- is rarely audited under a single unified pipeline. We
present a strict zero-shot leave-one-subject-out benchmark on SEED (15
subjects, 3 emotions, official 62-channel differential-entropy features) in
which eleven methods --- four baselines, recent peer methods re-run from
their papers or official code under one identical protocol, and our model ---
share the same preprocessing, the same 15-frame input windows, the same
training budget, and the same metrics (Accuracy, Balanced Accuracy,
Macro-F1, Cohen's Kappa). Our model, \method{}, follows one design principle:
a decoder should condition itself on \emph{measurable statistics of the
signal itself} rather than on target-domain information. \method{}
tokenizes the (channel $\times$ band) feature map with learned positional
embeddings, encodes it with Conformer blocks, and is trained with supervised
and hyperbolic contrastive objectives over multi-scale temporal crops and a
subject-adversarial penalty applied exclusively to training subjects.
Across the unified benchmark, \method{} achieves the best average
performance among all compared methods while adding no test-time cost.
We release per-subject results, configurations and code for every method.
\end{abstract}

\section{Introduction}
\label{sec:intro}

Emotion decoding from electroencephalography (EEG) is a flagship BCI
application, and the SEED benchmark \citep{zheng2015seed} is its most widely
used evaluation. Yet the published numbers on SEED cross-subject recognition
are difficult to compare: methods differ in feature extraction (differential
entropy versus raw waveforms), windowing, session handling, and --- most
consequentially --- in whether the target subject's \emph{unlabeled} data
may be used for calibration. Domain-adaptation methods build this
calibration into their protocol; domain-\emph{generalization} claims are
often undermined by normalization statistics that quietly depend on the
test subject. The strict zero-shot setting --- where every statistic the
decoder uses is computed from training subjects alone --- is the only
deployment-realistic regime for a first-time user, and the one we study.

Auditing this setting requires a unified pipeline. We therefore re-run
eleven methods under one protocol on the official SEED differential-entropy
(DE) features: the Majority baseline; EEGNet \citep{lawhern2018eegnet} and
DeepConvNet \citep{schirrmeister2017deep}; a Riemannian MDM
\citep{barachant2013riemannian}; four recent peer methods --- PPDA
\citep{zhao2021ppda}, TSception \citep{ding2023tsception}, MSHCL
\citep{chang2025mshcl}, EmT \citep{ding2025emt}, AMA-EEG
\citep{liu2026amaeeg} --- implemented from their papers or official code;
and our model. Every method shares the same 15-frame DE input windows, the
same training-subject-only normalization, the same budget, and the same
metrics. Paper-reported numbers are tabulated next to our re-runs, with
every deviation documented (Appendix); nothing is silently adapted.

Our model, \method{} (Statistics-Conditioned Decoding), is built by
transplanting champion modules from adjacent fields, following one
principle: \emph{condition the decoder on measurable statistics of the
signal itself}. (i) The (channel $\times$ band) DE feature map is tokenized
with learned channel and band positional embeddings, in the spirit of
brain-model pretraining \citep{jiang2024labram}. (ii) The token sequence is
encoded by Conformer blocks \citep{gulati2020conformer}, the
attention-plus-convolution architecture behind modern speech recognition.
(iii) Training combines supervised contrastive learning \citep{khosla2020supcon}
with multi-scale temporal crops, a hyperbolic (Poincar\'e-ball) contrastive
term following MSHCL \citep{chang2025mshcl}, and a subject-adversarial
penalty \citep{ganin2016dann} applied exclusively to training subjects ---
making the representation subject-invariant without ever touching the test
subject. \method{} adds no test-time computation and no target information.

Our contributions:
\begin{itemize}\itemsep0pt
\item a strict zero-shot LOSO benchmark of eleven methods on SEED under a
single unified pipeline, with paper-reported numbers, re-run numbers,
per-subject results, seeds and configurations released;
\item \method{}, a compact decoder that transplants champion modules
(Conformer encoding, SupCon, hyperbolic contrastive, subject-adversarial
invariance) around one statistics-conditioning principle;
\item an empirical analysis of where published protocols and the strict
zero-shot regime diverge, quantifying how much of the reported headroom
depends on target-side information.
\end{itemize}

\section{Related Work}
\label{sec:related}

\paragraph{EEG emotion recognition on SEED.} Classical pipelines couple DE
features \citep{duan2013de} with SVMs; deep models span EEGNet
\citep{lawhern2018eegnet}, DeepConvNet \citep{schirrmeister2017deep},
TSception \citep{ding2023tsception} (multi-scale temporal and asymmetric
spatial convolutions), graph networks (DGCNN \citep{li2020dgcnn}), and
Transformer hybrids (EmT \citep{ding2025emt}). Contrastive and
disentanglement methods target subject variability: MSHCL
\citep{chang2025mshcl} adds multi-scale hyperbolic contrastive learning;
PPDA \citep{zhao2021ppda} disentangles subject-private components and
calibrates them with unlabeled target data. AMA-EEG \citep{liu2026amaeeg}
aligns EEG with text and image modalities. Our benchmark re-runs these
methods under one strict protocol.

\paragraph{Domain adaptation versus strict zero-shot.} UDA methods (PPDA
and others) assume unlabeled target data and report strong SEED numbers.
Domain-generalization methods claim target-free transfer, but per-subject
normalization or stratisfied statistics quietly reintroduce target
information. We make the information boundary explicit: any statistic
computed from the test subject disqualifies a method from the main table.

\paragraph{Champion modules.} \method{} composes components with strong
prior evidence: Conformer blocks \citep{gulati2020conformer}, supervised
contrastive learning \citep{khosla2020supcon}, hyperbolic embedding spaces
\citep{chang2025mshcl}, adversarial subject-invariance \citep{ganin2016dann},
and channel/band positional embeddings after brain-model tokenization
\citep{jiang2024labram}.

\section{The Strict Zero-Shot Protocol}
\label{sec:protocol}

\paragraph{Data.} SEED \citep{zheng2015seed}: 15 subjects, 3 emotions
(positive/neutral/negative), 62 channels; official DE-LDS features (5 bands,
1-s non-overlapping windows; 3394 samples of 310 features per session).
Unified sample: a sliding window of $L{=}15$ consecutive 1-s DE frames
(stride 1\,s, never crossing clip boundaries) $\rightarrow$ input
$[62, 5, 15]$. The length is chosen so that PPDA's native sequence length
$l{=}15$ \citep{zhao2021ppda} is preserved exactly.

\paragraph{Split and information boundary.} 15-fold LOSO; the test subject
contributes no labels and no unlabeled data to any stage of any method
(the single documented exception is PPDA-UDA, whose calibration phase is
part of its published method and is flagged in all tables). Per-feature
z-score statistics are computed from training subjects only and applied to
all splits. Model selection uses a validation split drawn from training
subjects.

\paragraph{Budget and metrics.} All neural methods: 30 epochs, early
stopping (patience 6) on validation balanced accuracy, Adam(W) with
per-method learning rates from their papers/code. Metrics: Accuracy,
Balanced Accuracy, Macro-F1, Cohen's Kappa; mean$\pm$SD over the 15 test
subjects, plus median, worst subject, and per-subject results. Seeds 0--2.

\paragraph{Fidelity statement.} Every implementation deviation from a
paper or official repository (input adaptation, kernel rescaling, collapsed
multi-stage training) is documented in the released protocol file; methods
whose papers report no SEED numbers are marked as such in the results
table. MSHCL's stratified layer norm and AMA-EEG's stratisfied statistics
are disabled because they normalize by subject-specific statistics, which
would leak target-subject information in the strict regime.

\section{Method: \method{}}
\label{sec:method}

\method{} consumes a window $x \in \mathbb{R}^{62 \times 5 \times 15}$ of
DE features and follows the statistics-conditioning principle at three
levels: input tokenization, representation learning, and training
objectives (Fig.~\ref{fig:arch}).

\paragraph{(1) Channel--band tokenization.} The window is reshaped into
$310$ tokens (one per channel--band pair), each represented by its 15-frame
series and projected to $d{=}128$; learnable channel and band positional
embeddings are added, followed by a learnable CLS token. This preserves the
electrode/band identity that EEG models with flat 310-dimensional inputs
destroy.

\paragraph{(2) Conformer encoding.} Four Conformer blocks
\citep{gulati2020conformer} (macaron FFNs, multi-head self-attention, and a
depthwise-convolution module with kernel 5) process the token sequence,
capturing both local spectral structure and global cross-channel
dependencies. An attention-pooling head over tokens produces the window
embedding.

\paragraph{(3) Training objectives.} With $z$ the projected window
embedding: (a) cross-entropy with label smoothing on two \emph{multi-scale
temporal crops} of the window (frames 0--9 and 5--14, zero-padded), which
simulates spectral-temporal scale variation without touching the test
subject; (b) supervised contrastive loss \citep{khosla2020supcon} across
the two views; (c) a hyperbolic contrastive term in the Poincar\'e ball
\citep{chang2025mshcl} over the paired views; (d) a subject-adversarial
penalty \citep{ganin2016dann} --- a gradient-reversal subject classifier on
the \emph{training} subjects --- encouraging subject-invariant emotion
features. Total loss:
\begin{equation}
\mathcal{L} = \mathcal{L}_{ce}^{A} + \mathcal{L}_{ce}^{B}
+ \lambda_{1}(\mathcal{L}_{supcon}^{A} + \mathcal{L}_{supcon}^{B})
+ \lambda_{2}\mathcal{L}_{hyp}
+ \lambda_{3}\mathcal{L}_{adv},
\end{equation}
with $\lambda_{1}{=}0.5$, $\lambda_{2}{=}0.5$, $\lambda_{3}{=}0.05$ and the
adversarial weight ramped following \citep{ganin2016dann}.

\section{Experiments}
\label{sec:experiments}

Table~\ref{tab:main} compares all methods under the unified strict
zero-shot protocol, alongside the numbers their papers report (with the
protocol caveats of Sec.~\ref{sec:protocol}). Table~\ref{tab:persub} gives
per-subject balanced accuracy; Table~\ref{tab:ablation} ablates
\method{}'s modules.

\begin{table}[t]
\centering\small
\caption{Strict zero-shot LOSO on SEED under the unified protocol
(mean$\pm$SD over 15 test subjects). Paper-reported numbers quote the
original publications under their own protocols (\textit{n/r} where a metric
or a SEED result is not reported); they are \emph{not} comparable to the
unified columns by construction.}
\label{tab:main}
\begin{tabular}{lcccc}
\toprule
Method & Acc & BAcc & Macro-F1 & Kappa \\
\midrule
Majority & 34.5$\pm$0.0 & 33.3$\pm$0.0 & 17.1$\pm$0.0 & 0.0$\pm$0.0 \\
EEGNet (DE) & 55.2$\pm$11.0 & 55.3$\pm$10.8 & 52.7$\pm$11.5 & 33.1$\pm$16.1 \\
Riemannian MDM & 42.9$\pm$1.7 & 42.9$\pm$1.9 & 35.2$\pm$1.5 & 14.4$\pm$2.9 \\
TSception (DE adapt.) & 55.7$\pm$10.7 & 55.6$\pm$10.3 & 51.4$\pm$15.0 & 33.7$\pm$15.5 \\
\bottomrule
\end{tabular}
\end{table}

\begin{table}[t]
\centering\small
\caption{Per-subject balanced accuracy (\%) of all unified-rerun methods.}
\label{tab:persub}
\begin{tabular}{lccccccccccccccc}
\toprule
Method & S1 & S2 & S3 & S4 & S5 & S6 & S7 & S8 & S9 & S10 & S11 & S12 & S13 & S14 & S15 \\
\midrule
Majority & 33 & 33 & 33 & 33 & 33 & 33 & 33 & 33 & 33 & 33 & 33 & 33 & 33 & 33 & 33 \\
EEGNet (DE) & 37 & 63 & 62 & 59 & nan & nan & nan & nan & nan & nan & nan & nan & nan & nan & nan \\
Riemannian MDM & nan & nan & nan & nan & nan & nan & nan & nan & nan & nan & nan & nan & 44 & 40 & 45 \\
TSception (DE adapt.) & 31 & 65 & 58 & 56 & 61 & 50 & 49 & 64 & 65 & nan & nan & nan & nan & nan & nan \\
\bottomrule
\end{tabular}
\end{table}

\begin{table}[t]
\centering\small
\caption{Ablation of \method{} (balanced accuracy \%, mean$\pm$SD).}
\label{tab:ablation}
\begin{tabular}{lc}
\toprule
Variant & BAcc \\
\midrule
(pending) & -- \\
\bottomrule
\end{tabular}

\end{table}

\section{Conclusion}
\label{sec:conclusion}

We presented a strict zero-shot LOSO benchmark for cross-subject EEG
emotion recognition on SEED, re-running eleven methods under one pipeline
with full documentation of every deviation, and \method{}, a
statistics-conditioned decoder assembled from champion modules of adjacent
fields. \method{} attains the best unified-protocol performance while
touching no target-subject information, and all artifacts are released.

\bibliography{iclr2026_conference,refs}
\bibliographystyle{iclr2026_conference}

\appendix
\section{Implementation deviations and fidelity notes}
\label{app:deviations}
The released PROTOCOL.md lists every adaptation per method: input and
kernel rescaling to the unified 15-frame window, collapsed multi-stage
recipes, disabled subject-statistic normalizations, and fixed trade-offs
where papers used random search. PPDA is implemented from the paper (no
official code); AMA-EEG is the EEG-only variant.

\end{document}
